\documentclass[12pt,a4paper]{article}
\usepackage[UTF8]{ctex}
\usepackage{geometry}
\usepackage{listings}
\usepackage{xcolor}
\usepackage{hyperref}
\usepackage{booktabs}
\usepackage{float}
\usepackage{tcolorbox}
\usepackage{enumitem}
\usepackage{graphicx}
\usepackage{longtable}

\geometry{left=2.5cm,right=2.5cm,top=2.5cm,bottom=2.5cm}

\lstset{
    language=SQL,
    basicstyle=\ttfamily\small,
    keywordstyle=\color{blue}\bfseries,
    commentstyle=\color{gray},
    stringstyle=\color{orange},
    numbers=left,
    numberstyle=\tiny\color{gray},
    stepnumber=1,
    numbersep=8pt,
    backgroundcolor=\color{gray!10},
    frame=single,
    frameround=fttt,
    breaklines=true,
    showstringspaces=false,
    tabsize=4,
    captionpos=b,
    xleftmargin=2em,
    xrightmargin=1em,
    aboveskip=1em,
    belowskip=1em
}

\tcbuselibrary{skins,breakable}
\newtcolorbox{knowledgebox}[1][]{
    colback=blue!5!white, colframe=blue!50!black,
    fonttitle=\bfseries, title=#1, breakable}
\newtcolorbox{practicebox}[1][]{
    colback=green!5!white, colframe=green!50!black,
    fonttitle=\bfseries, title=#1, breakable}
\newtcolorbox{tipbox}[1][]{
    colback=yellow!10!white, colframe=orange!70!black,
    fonttitle=\bfseries, title=#1, breakable}
\newtcolorbox{theorybox}[1][]{
    colback=purple!5!white, colframe=purple!70!black,
    fonttitle=\bfseries, title=#1, breakable}
\newtcolorbox{warnbox}[1][]{
    colback=red!5!white, colframe=red!70!black,
    fonttitle=\bfseries, title=#1, breakable}

\hypersetup{colorlinks=true, linkcolor=blue, urlcolor=cyan, bookmarks=true}

\title{\textbf{大数据分析实战手册}\\[0.5em]\large 从数据探索到智能分析}
\author{竞赛培训组}
\date{\today}

\begin{document}
\maketitle

\begin{tipbox}[学习指引]
\textbf{前序课程}：确保你已经完成了大数据平台的部署（见《大数据平台部署实战手册》）。

\textbf{后续课程}：分析完成后，下一步是将结果可视化。

\textbf{配套练习}：更多练习见仓库根目录的《仿真实战练习.md》。
\end{tipbox}

\tableofcontents
\newpage

%========================================
\section{前言：什么是大数据分析？}
%========================================

\begin{theorybox}[用体检报告理解大数据分析]
你去医院体检，医生给你做了一系列检查（血压、血糖、心率...）。这些检查结果就是"数据"。

医生做三件事：
\begin{enumerate}
    \item \textbf{看数据}：血压多少、血糖多少（数据探索）
    \item \textbf{找规律}：血压偏高的人通常体重也偏高（关联分析）
    \item \textbf{下结论}：建议你少吃甜食、多运动（分析结论）
\end{enumerate}

\textbf{大数据分析}就是同样的过程——用计算机代替医生，对海量数据进行"看数据→找规律→下结论"。
\end{theorybox}

\begin{knowledgebox}[竞赛中的数据分析考什么？]
根据竞赛评分标准，数据分析占\textbf{30分}。要求你：
\begin{itemize}
    \item 在已部署的大数据平台上，用 SQL 查询数据
    \item 对数据进行多维度统计分析
    \item 发现数据中的规律和趋势
    \item 得出有价值的分析结论
\end{itemize}

\textbf{你需要掌握的工具}：
\begin{itemize}
    \item \textbf{Hive SQL}：在 HDFS 上用 SQL 查询数据（适合离线大批量分析）
    \item \textbf{Spark SQL}：用 Spark 引擎执行 SQL（比 Hive 快，适合中等规模分析）
    \item \textbf{PySpark}：用 Python 调用 Spark（适合做数据处理+机器学习）
\end{itemize}
\end{knowledgebox}

\begin{tipbox}[分析流程一图看懂]
\begin{footnotesize}
\begin{verbatim}
[HDFS数据] -> [Hive建表] -> [SQL分析] -> [导出] -> [报告]
  原材料       标签系统     医生诊断     报告     论文
\end{verbatim}
\end{footnotesize}
\end{tipbox}

%========================================
\section{第一章：数据准备}
%========================================

\begin{knowledgebox}[这一章要做什么？]
数据分析的前提是"数据已经存好了"。这一章教你：
\begin{enumerate}
    \item 从云端仿真平台获取数据
    \item 把数据上传到 HDFS（分布式文件系统）
    \item 在 Hive 中创建表格（让数据可以用 SQL 查询）
    \item 验证数据是否正确
\end{enumerate}
\end{knowledgebox}

\subsection{步骤1：从仿真平台获取数据}

我们的数据来自云端电商数据仿真平台，平台每秒自动生成用户行为数据。

\begin{lstlisting}[language=bash]
# 【在 node1 上执行】从仿真平台获取数据
# 获取 1000 条数据，保存为本地 JSON 文件
curl -s "http://163.7.1.176:8889/api/data?n=1000" > /tmp/user_behavior.json

# 查看前几条数据
head -c 500 /tmp/user_behavior.json
\end{lstlisting}

\subsection{步骤2：将数据上传到 HDFS}

\begin{theorybox}[为什么要上传到 HDFS？]
HDFS 是分布式文件系统，数据存到 HDFS 上后，所有节点都能访问。Hive 和 Spark 也都是从 HDFS 上读取数据。

\textbf{类比}：HDFS 是图书馆的书架，Hive 是图书馆的目录卡片，Spark 是图书馆的分析员。
\end{theorybox}

\begin{lstlisting}[language=bash]
# 【在 node1 上执行】创建 HDFS 目录
hdfs dfs -mkdir -p /user/data/ecommerce

# 上传数据到 HDFS
hdfs dfs -put /tmp/user_behavior.json /user/data/ecommerce/

# 验证上传成功
hdfs dfs -ls /user/data/ecommerce/
# 应该看到 user_behavior.json 文件

# 查看文件前几行
hdfs dfs -cat /user/data/ecommerce/user_behavior.json | head -5
\end{lstlisting}

\subsection{步骤3：在 Hive 中创建表}

\begin{knowledgebox}[Hive 建表是什么意思？]
HDFS 上的文件是"裸数据"（纯文本），不能直接用 SQL 查询。

Hive 建表就是给这些数据"贴标签"——告诉 Hive：
\begin{itemize}
    \item 第1列叫 timestamp，是字符串类型
    \item 第2列叫 user\_id，是字符串类型
    \item 第3列叫 event\_type，是字符串类型
    \item ...依此类推
\end{itemize}

建表之后，你就可以写 SQL 查询了：\texttt{SELECT * FROM user\_behavior WHERE event\_type='buy'}
\end{knowledgebox}

\begin{lstlisting}[language=SQL]
-- 【在 node1 上执行】启动 Hive
hive

-- 创建数据库
CREATE DATABASE IF NOT EXISTS ecommerce;
USE ecommerce;

-- 创建外部表（数据在 HDFS 上，表只是"标签"）
CREATE EXTERNAL TABLE user_behavior (
    user_id STRING,        -- 用户ID
    event_type STRING,     -- 事件类型：view/click/add_cart/buy/search/collect/share
    session_id STRING,     -- 会话ID
    device STRING,         -- 设备：iOS/Android/PC/小程序/H5
    ip STRING,             -- IP地址
    product_id STRING,     -- 商品ID
    category STRING,       -- 商品类别
    price DOUBLE,          -- 商品价格
    page_source STRING,    -- 来源页面
    stay_time INT          -- 停留秒数
)
ROW FORMAT SERDE
    'org.apache.hive.hcatalog.data.JsonSerDe'
LOCATION '/user/data/ecommerce/';
\end{lstlisting}

\begin{warnbox}[常见错误]
\begin{itemize}
    \item 报错 \texttt{JsonSerDe not found}：需要在 Hive 中添加 JSON SerDe jar 包
    \item 解决办法：\texttt{ADD JAR \$HIVE\_HOME/hcatalog/.../hive-hcatalog-core-*.jar;}
\end{itemize}
\end{warnbox}

\subsection{步骤4：验证数据}

\begin{lstlisting}[language=SQL]
-- 查看表结构
DESCRIBE user_behavior;

-- 查看前10条数据
SELECT * FROM user_behavior LIMIT 10;

-- 统计总记录数
SELECT COUNT(*) FROM user_behavior;

-- 查看有哪些事件类型
SELECT DISTINCT event_type FROM user_behavior;

-- 查看有哪些商品类别
SELECT DISTINCT category FROM user_behavior;
\end{lstlisting}

\begin{practicebox}[动手练习]
\begin{enumerate}
    \item 从仿真平台获取 5000 条数据
    \item 上传到 HDFS
    \item 在 Hive 中建表并验证
    \item 执行以上 5 个 SQL 查询，观察结果
\end{enumerate}
\end{practicebox}

%========================================
\section{第二章：数据探索——SQL 基础查询}
%========================================

\begin{knowledgebox}[这一章要做什么？]
学会用 SQL 对数据进行基本的统计分析。这是数据分析的基本功，就像医生要学会看体检报告一样。
\end{knowledgebox}

\subsection{基本统计函数}

\begin{lstlisting}[language=SQL]
-- 统计总记录数
SELECT COUNT(*) AS 总记录数 FROM user_behavior;

-- 统计去重后的用户数
SELECT COUNT(DISTINCT user_id) AS 独立用户数 FROM user_behavior;

-- 统计去重后的商品数
SELECT COUNT(DISTINCT product_id) AS 独立商品数 FROM user_behavior;

-- 价格统计（平均值、最大值、最小值）
SELECT
    AVG(price) AS 平均价格,
    MAX(price) AS 最高价格,
    MIN(price) AS 最低价格,
    STDDEV(price) AS 价格标准差
FROM user_behavior
WHERE price IS NOT NULL;

-- 停留时间统计
SELECT
    AVG(stay_time) AS 平均停留秒数,
    MAX(stay_time) AS 最长停留秒数,
    MIN(stay_time) AS 最短停留秒数
FROM user_behavior
WHERE stay_time IS NOT NULL;
\end{lstlisting}

\subsection{分组统计（GROUP BY）}

\begin{theorybox}[GROUP BY 是什么？]
\texttt{GROUP BY} = "按某个维度分组，然后对每组做统计"。

\textbf{类比}：老师把全班同学按"性别"分组，然后统计每组的平均成绩。这就是 GROUP BY。
\end{theorybox}

\begin{lstlisting}[language=SQL]
-- 按事件类型分组，统计每种事件的数量
SELECT
    event_type AS 事件类型,
    COUNT(*) AS 数量
FROM user_behavior
GROUP BY event_type
ORDER BY 数量 DESC;

-- 按商品类别分组，统计每类的平均价格
SELECT
    category AS 品类,
    COUNT(*) AS 记录数,
    ROUND(AVG(price), 2) AS 平均价格
FROM user_behavior
WHERE price IS NOT NULL
GROUP BY category
ORDER BY 记录数 DESC;

-- 按设备类型分组
SELECT
    device AS 设备,
    COUNT(*) AS 使用次数,
    ROUND(COUNT(*) * 100.0 / (SELECT COUNT(*) FROM user_behavior), 2) AS 占比百分比
FROM user_behavior
GROUP BY device
ORDER BY 使用次数 DESC;
\end{lstlisting}

\subsection{条件过滤（WHERE / HAVING）}

\begin{lstlisting}[language=SQL]
-- WHERE：过滤行（分组前过滤）
-- 查看所有"购买"事件
SELECT * FROM user_behavior WHERE event_type = 'buy';

-- 查看价格大于 5000 的记录
SELECT * FROM user_behavior WHERE price > 5000;

-- 查看 iOS 用户的行为
SELECT * FROM user_behavior WHERE device = 'iOS';

-- 多条件过滤（AND / OR）
SELECT * FROM user_behavior
WHERE event_type = 'buy' AND price > 1000;

-- HAVING：过滤组（分组后过滤）
-- 只显示记录数大于 100 的品类
SELECT
    category,
    COUNT(*) AS 数量
FROM user_behavior
GROUP BY category
HAVING COUNT(*) > 100
ORDER BY 数量 DESC;
\end{lstlisting}

\subsection{排序与限制}

\begin{lstlisting}[language=SQL]
-- ORDER BY：排序
-- 按价格降序排列（从高到低）
SELECT * FROM user_behavior
WHERE price IS NOT NULL
ORDER BY price DESC
LIMIT 20;

-- 按停留时间升序排列（从低到高）
SELECT * FROM user_behavior
WHERE stay_time IS NOT NULL
ORDER BY stay_time ASC
LIMIT 20;
\end{lstlisting}

%========================================
\section{第三章：数据分析实战}
%========================================

\begin{knowledgebox}[这一章要做什么？]
用真实的电商数据，做 5 个具体的分析场景。每个场景都包含：业务问题 $\rightarrow$ SQL 查询 $\rightarrow$ 结果解读。

这是竞赛中最可能考到的内容！
\end{knowledgebox}

\subsection{场景1：用户行为分析}

\begin{theorybox}[业务问题]
老板问："我们的用户在网站上都在做什么？哪些行为最多？"
\end{theorybox}

\begin{lstlisting}[language=SQL]
-- 分析1：各种行为的占比
SELECT
    event_type AS 行为类型,
    COUNT(*) AS 次数,
    ROUND(COUNT(*) * 100.0 / (SELECT COUNT(*) FROM user_behavior), 2) AS 占比
FROM user_behavior
GROUP BY event_type
ORDER BY 次数 DESC;

-- 分析2：每种行为的平均价格
SELECT
    event_type AS 行为类型,
    ROUND(AVG(price), 2) AS 平均价格,
    ROUND(MIN(price), 2) AS 最低价格,
    ROUND(MAX(price), 2) AS 最高价格
FROM user_behavior
WHERE price IS NOT NULL
GROUP BY event_type
ORDER BY 平均价格 DESC;

-- 分析3：各品类的行为分布
SELECT
    category AS 品类,
    event_type AS 行为,
    COUNT(*) AS 次数
FROM user_behavior
GROUP BY category, event_type
ORDER BY category, 次数 DESC;
\end{lstlisting}

\subsection{场景2：用户画像分析}

\begin{theorybox}[业务问题]
老板问："我们的用户都是什么人？用什么设备？在哪些城市？"
\end{theorybox}

\begin{lstlisting}[language=SQL]
-- 分析1：设备分布
SELECT
    device AS 设备类型,
    COUNT(*) AS 使用次数,
    ROUND(COUNT(*) * 100.0 / (SELECT COUNT(*) FROM user_behavior), 2) AS 占比
FROM user_behavior
GROUP BY device
ORDER BY 使用次数 DESC;

-- 分析2：来源页面分布
SELECT
    page_source AS 来源页面,
    COUNT(*) AS 访问次数,
    ROUND(COUNT(*) * 100.0 / (SELECT COUNT(*) FROM user_behavior), 2) AS 占比
FROM user_behavior
GROUP BY page_source
ORDER BY 访问次数 DESC;

-- 分析3：品类偏好（用户最关注哪些品类）
SELECT
    category AS 品类,
    COUNT(DISTINCT user_id) AS 独立用户数,
    COUNT(*) AS 总行为数
FROM user_behavior
GROUP BY category
ORDER BY 独立用户数 DESC;
\end{lstlisting}

\subsection{场景3：商品分析}

\begin{theorybox}[业务问题]
老板问："哪些商品最受欢迎？哪个品类卖得最好？"
\end{theorybox}

\begin{lstlisting}[language=SQL]
-- 分析1：热销商品 TOP 10
SELECT
    product_id AS 商品ID,
    category AS 品类,
    COUNT(*) AS 被操作次数,
    SUM(CASE WHEN event_type = 'buy' THEN 1 ELSE 0 END) AS 购买次数,
    SUM(CASE WHEN event_type = 'view' THEN 1 ELSE 0 END) AS 浏览次数
FROM user_behavior
GROUP BY product_id, category
ORDER BY 购买次数 DESC, 被操作次数 DESC
LIMIT 10;

-- 分析2：品类销量排行
SELECT
    category AS 品类,
    SUM(CASE WHEN event_type = 'buy' THEN 1 ELSE 0 END) AS 购买次数,
    SUM(CASE WHEN event_type = 'view' THEN 1 ELSE 0 END) AS 浏览次数,
    ROUND(SUM(CASE WHEN event_type = 'buy' THEN 1 ELSE 0 END) * 100.0 /
          SUM(CASE WHEN event_type = 'view' THEN 1 ELSE 0 END), 2) AS 转化率
FROM user_behavior
GROUP BY category
ORDER BY 购买次数 DESC;

-- 分析3：价格区间分布
SELECT
    CASE
        WHEN price < 100 THEN '0-100元'
        WHEN price < 500 THEN '100-500元'
        WHEN price < 1000 THEN '500-1000元'
        WHEN price < 5000 THEN '1000-5000元'
        ELSE '5000元以上'
    END AS 价格区间,
    COUNT(*) AS 商品数量,
    ROUND(AVG(price), 2) AS 平均价格
FROM user_behavior
WHERE price IS NOT NULL
GROUP BY
    CASE
        WHEN price < 100 THEN '0-100元'
        WHEN price < 500 THEN '100-500元'
        WHEN price < 1000 THEN '500-1000元'
        WHEN price < 5000 THEN '1000-5000元'
        ELSE '5000元以上'
    END
ORDER BY 平均价格;
\end{lstlisting}

\subsection{场景4：转化率分析（漏斗分析）}

\begin{theorybox}[业务问题]
老板问："100个浏览的人，有多少人加了购物车？有多少人最终买了？"

这就是"漏斗分析"——从浏览 $\rightarrow$ 加购 $\rightarrow$ 购买，每一步有多少人流失。
\end{theorybox}

\begin{lstlisting}[language=SQL]
-- 漏斗分析：浏览 -> 加购 -> 购买
SELECT
    '1.浏览' AS 步骤,
    COUNT(DISTINCT user_id) AS 用户数
FROM user_behavior WHERE event_type = 'view'
UNION ALL
SELECT
    '2.加购' AS 步骤,
    COUNT(DISTINCT user_id) AS 用户数
FROM user_behavior WHERE event_type = 'add_cart'
UNION ALL
SELECT
    '3.购买' AS 步骤,
    COUNT(DISTINCT user_id) AS 用户数
FROM user_behavior WHERE event_type = 'buy';

-- 各品类转化率
SELECT
    category AS 品类,
    COUNT(DISTINCT CASE WHEN event_type='view' THEN user_id END) AS 浏览用户,
    COUNT(DISTINCT CASE WHEN event_type='add_cart' THEN user_id END) AS 加购用户,
    COUNT(DISTINCT CASE WHEN event_type='buy' THEN user_id END) AS 购买用户,
    ROUND(COUNT(DISTINCT CASE WHEN event_type='buy' THEN user_id END) * 100.0 /
          NULLIF(COUNT(DISTINCT CASE WHEN event_type='view' THEN user_id END), 0), 2)
          AS 浏览到购买转化率
FROM user_behavior
GROUP BY category
ORDER BY 浏览到购买转化率 DESC;
\end{lstlisting}

\subsection{场景5：时间趋势分析}

\begin{theorybox}[业务问题]
老板问："用户一般什么时候最活跃？上午还是下午？"
\end{theorybox}

\begin{lstlisting}[language=SQL]
-- 按小时统计活跃度（需要先从 timestamp 提取小时）
SELECT
    SUBSTR(timestamp, 12, 2) AS 小时,
    COUNT(*) AS 行为次数,
    COUNT(DISTINCT user_id) AS 独立用户数
FROM user_behavior
GROUP BY SUBSTR(timestamp, 12, 2)
ORDER BY 小时;

-- 按日期统计趋势
SELECT
    SUBSTR(timestamp, 1, 10) AS 日期,
    COUNT(*) AS 行为次数,
    COUNT(DISTINCT user_id) AS 独立用户数,
    SUM(CASE WHEN event_type = 'buy' THEN 1 ELSE 0 END) AS 购买次数
FROM user_behavior
GROUP BY SUBSTR(timestamp, 1, 10)
ORDER BY 日期;
\end{lstlisting}

%========================================
\section{第四章：简单机器学习（Spark MLlib）}
%========================================

\begin{knowledgebox}[这一章要做什么？]
用 Spark 的机器学习库（MLlib），对用户数据做一个简单的预测模型。比如"预测用户是否会购买"。

\textbf{类比}：医生根据体检数据（特征），预测你是否会得某种病（标签）。这里我们根据用户行为数据（特征），预测用户是否会购买（标签）。
\end{knowledgebox}

\subsection{数据预处理}

\begin{lstlisting}[language=Python]
# 【在 node1 上执行】启动 pyspark
pyspark

from pyspark.sql import SparkSession
from pyspark.sql.functions import *

# 读取 HDFS 上的数据
spark = SparkSession.builder.appName("UserAnalysis").getOrCreate()
df = spark.read.json("hdfs://node1:9000/user/data/ecommerce/user_behavior.json")

# 查看数据
df.show(5)
df.printSchema()

# 创建特征：每个用户的行为统计
user_features = df.groupBy("user_id").agg(
    count("*").alias("总行为数"),
    sum(when(col("event_type") == "view", 1).otherwise(0)).alias("浏览次数"),
    sum(when(col("event_type") == "click", 1).otherwise(0)).alias("点击次数"),
    sum(when(col("event_type") == "add_cart", 1).otherwise(0)).alias("加购次数"),
    sum(when(col("event_type") == "buy", 1).otherwise(0)).alias("购买次数"),
    avg("price").alias("平均价格"),
    avg("stay_time").alias("平均停留时间")
)

# 创建标签：是否购买过（1=是，0=否）
user_features = user_features.withColumn(
    "是否购买", when(col("购买次数") > 0, 1).otherwise(0)
)

user_features.show(10)
\end{lstlisting}

\subsection{训练分类模型}

\begin{lstlisting}[language=Python]
from pyspark.ml.feature import VectorAssembler
from pyspark.ml.classification import LogisticRegression
from pyspark.ml.evaluation import BinaryClassificationEvaluator

# 1. 组装特征向量（MLlib 要求特征放在一个向量列中）
assembler = VectorAssembler(
    inputCols=["总行为数", "浏览次数", "点击次数", "加购次数", "平均价格", "平均停留时间"],
    outputCol="features"
)
data = assembler.transform(user_features)

# 2. 拆分训练集和测试集（70%训练，30%测试）
train, test = data.randomSplit([0.7, 0.3], seed=42)
print(f"训练集: {train.count()} 条, 测试集: {test.count()} 条")

# 3. 训练模型（逻辑回归）
lr = LogisticRegression(featuresCol="features", labelCol="是否购买")
model = lr.fit(train)

# 4. 在测试集上预测
predictions = model.transform(test)

# 5. 评估模型准确率
evaluator = BinaryClassificationEvaluator(labelCol="是否购买")
auc = evaluator.evaluate(predictions)
print(f"模型 AUC: {auc:.4f}")

# 查看预测结果
predictions.select("user_id", "是否购买", "prediction", "probability").show(10)
\end{lstlisting}

\subsection{结果解读}

\begin{tipbox}[AUC 是什么意思？]
AUC（Area Under Curve）是分类模型的评价指标，取值 0-1：
\begin{itemize}
    \item AUC = 0.5：模型和瞎猜一样（没用）
    \item AUC = 0.7：还可以
    \item AUC = 0.8：不错
    \item AUC = 0.9+：很好
\end{itemize}
\end{tipbox}

%========================================
\section{第五章：数据导出与报告}
%========================================

\subsection{导出分析结果}

\begin{lstlisting}[language=bash]
# 【在 node1 上执行】将 Hive 查询结果导出为 CSV
hive -e "
SET hive.resultset.use.unique.column.names=false;
SELECT event_type, COUNT(*) as cnt
FROM ecommerce.user_behavior
GROUP BY event_type
ORDER BY cnt DESC
" | sed 's/\t/,/g' > /tmp/event_stats.csv

# 上传到 HDFS 或下载到本地
cat /tmp/event_stats.csv
\end{lstlisting}

\subsection{用 Python 做可视化}

\begin{lstlisting}[language=Python]
# 【在本地电脑上执行】用 matplotlib 画图
import matplotlib.pyplot as plt
import pandas as pd

# 读取分析结果
df = pd.read_csv('/tmp/event_stats.csv')

# 柱状图
plt.figure(figsize=(10, 6))
plt.bar(df['event_type'], df['cnt'], color='steelblue')
plt.title('用户行为分布')
plt.xlabel('行为类型')
plt.ylabel('次数')
plt.xticks(rotation=45)
plt.tight_layout()
plt.savefig('event_distribution.png', dpi=150)
plt.show()
\end{lstlisting}

\subsection{竞赛报告框架}

\begin{knowledgebox}[报告应该包含什么？]
竞赛报告是评分的重要部分（占 30 分）。一份好的报告应包含：
\begin{enumerate}
    \item \textbf{问题描述}：你分析了什么问题
    \item \textbf{数据说明}：数据来源、字段含义、数据量
    \item \textbf{分析方法}：用了什么工具、什么 SQL 查询
    \item \textbf{分析结果}：图表 + 数字 + 文字说明
    \item \textbf{结论建议}：根据分析结果，给出什么建议
\end{enumerate}
\end{knowledgebox}

%========================================
\section{第六章：从不同数据源读取数据}

\begin{knowledgebox}[为什么要学多种数据源？]
大数据平台上的数据可能存储在不同地方：
\begin{itemize}
    \item \textbf{HDFS}：原始文件（CSV、JSON、日志）
    \item \textbf{HBase}：结构化的用户/商品信息
    \item \textbf{ClickHouse}：已经统计好的聚合数据
\end{itemize}
你需要知道"从哪里读、怎么读"。
\end{knowledgebox}

\subsection{从 HDFS 读取（通过 Hive）}

\begin{lstlisting}[language=SQL]
-- Hive 通过 LOCATION 指向 HDFS 上的文件
-- 建表时指定 HDFS 路径，之后就可以用 SQL 查询
CREATE EXTERNAL TABLE raw_logs (
    line STRING
) LOCATION '/user/data/logs/';

-- 也可以用正则提取字段
CREATE EXTERNAL TABLE parsed_logs (
    user_id STRING, action STRING, time_str STRING
)
ROW FORMAT SERDE 'org.apache.hadoop.hive.serde2.RegexSerDe'
WITH SERDEPROPERTIES (
    "input.regex" = "(\\S+)\\s+(\\S+)\\s+(\\S+)"
)
LOCATION '/user/data/logs/';
\end{lstlisting}

\subsection{从 HBase 读取}

\begin{lstlisting}[language=SQL]
-- 在 Hive 中创建外部表映射 HBase
CREATE EXTERNAL TABLE user_profile_hbase (
    rowkey STRING,
    name STRING,
    age INT,
    city STRING
)
STORED BY 'org.apache.hadoop.hive.hbase.HBaseStorageHandler'
WITH SERDEPROPERTIES (
    "hbase.columns.mapping" = ":key,info:name,info:age,info:city"
)
TBLPROPERTIES ("hbase.table.name" = "user_profile");

-- 现在可以用 SQL 查询 HBase 数据了
SELECT city, AVG(age) FROM user_profile_hbase GROUP BY city;
\end{lstlisting}

\subsection{PySpark 读取 HDFS JSON}

\begin{lstlisting}[language=Python]
from pyspark.sql import SparkSession

spark = SparkSession.builder.appName("Analysis").getOrCreate()

# 读取 HDFS 上的 JSON 文件
df = spark.read.json("hdfs://node1:9000/user/data/ecommerce/user_behavior.json")

# 查看数据
df.show(5)           # 前5行
df.printSchema()     # 列名和类型
df.describe().show() # 统计摘要
df.count()           # 总行数
\end{lstlisting}

\section{第七章：Spark DataFrame 操作}

\begin{knowledgebox}[Spark DataFrame 是什么？]
Spark DataFrame = "分布式版本的 Pandas DataFrame"。语法非常相似，但数据分布在多台机器上并行处理。

\textbf{类比}：Pandas 是一个工人干活，Spark 是一群工人同时干活。
\end{knowledgebox}

\subsection{基本操作}

\begin{lstlisting}[language=Python]
from pyspark.sql import SparkSession
from pyspark.sql.functions import *

spark = SparkSession.builder.getOrCreate()
df = spark.read.json("hdfs://node1:9000/user/data/ecommerce/user_behavior.json")

# 选择列
df.select("user_id", "event_type", "price").show(5)

# 过滤行
df.filter(col("event_type") == "buy").show(5)
df.filter(col("price") > 1000).show(5)
df.filter((col("event_type") == "buy") & (col("price") > 1000)).show(5)

# 分组统计
df.groupBy("event_type").count().show()
df.groupBy("category").agg(
    avg("price").alias("平均价格"),
    count("*").alias("记录数")
).show()

# 排序
df.orderBy(col("price").desc()).show(10)
\end{lstlisting}

\subsection{数据转换}

\begin{lstlisting}[language=Python]
# 新增列
df = df.withColumn("price_level",
    when(col("price") < 100, "低价")
    .when(col("price") < 1000, "中价")
    .otherwise("高价")
)

# 删除列
df = df.drop("ip")

# 重命名列
df = df.withColumnRenamed("event_type", "事件类型")

# 类型转换
df = df.withColumn("price", col("price").cast("double"))
\end{lstlisting}

\subsection{多表关联（JOIN）}

\begin{lstlisting}[language=Python]
# 假设有两个表：用户行为表 和 商品信息表
behavior = spark.read.json("hdfs://node1:9000/user/data/ecommerce/user_behavior.json")
# products 需要先从 HDFS 读取（或从 Hive 查询）

# 内连接：只保留两个表都匹配的行
joined = behavior.join(products, on="product_id", how="inner")

# 左连接：保留左表所有行
joined = behavior.join(products, on="product_id", how="left")

# 关联后统计
joined.groupBy("category").agg(
    count("*").alias("总行为数"),
    sum(when(col("event_type")=="buy", 1).otherwise(0)).alias("购买数")
).show()
\end{lstlisting}

\section{第八章：更多分析场景}

\subsection{场景6：RFM 用户价值分析}

\begin{theorybox}[什么是 RFM？]
RFM 是衡量用户价值的经典模型：
\begin{itemize}
    \item \textbf{R}（Recency）：最近一次消费距今多久（越近越好）
    \item \textbf{F}（Frequency）：消费频率（越多越好）
    \item \textbf{M}（Monetary）：消费金额（越高越好）
\end{itemize}
\end{theorybox}

\begin{lstlisting}[language=SQL]
-- 在 Hive 中计算 RFM
SELECT
    user_id,
    MAX(timestamp) AS 最近消费时间,
    COUNT(CASE WHEN event_type='buy' THEN 1 END) AS 消费次数,
    SUM(CASE WHEN event_type='buy' THEN price ELSE 0 END) AS 消费总额
FROM user_behavior
GROUP BY user_id
ORDER BY 消费总额 DESC
LIMIT 20;
\end{lstlisting}

\subsection{场景7：留存分析}

\begin{theorybox}[什么是留存？]
留存 = "昨天来过的用户，今天还来吗？"。留存率越高，说明用户粘性越好。
\end{theorybox}

\begin{lstlisting}[language=SQL]
-- 按日期统计独立用户数和新用户数
SELECT
    SUBSTR(timestamp, 1, 10) AS 日期,
    COUNT(DISTINCT user_id) AS 独立用户数
FROM user_behavior
GROUP BY SUBSTR(timestamp, 1, 10)
ORDER BY 日期;
\end{lstlisting}

\subsection{场景8：品类关联分析}

\begin{lstlisting}[language=SQL]
-- 找出被同一批用户关注最多的品类对
SELECT
    a.category AS 品类A,
    b.category AS 品类B,
    COUNT(DISTINCT a.user_id) AS 共同用户数
FROM user_behavior a
JOIN user_behavior b
    ON a.user_id = b.user_id
    AND a.category < b.category
WHERE a.event_type = 'buy' AND b.event_type = 'buy'
GROUP BY a.category, b.category
ORDER BY 共同用户数 DESC
LIMIT 20;
\end{lstlisting}

\section{第九章：竞赛报告写作指导}

\begin{knowledgebox}[报告结构模板（5 部分）]
\begin{enumerate}
    \item \textbf{问题描述}（200字）：我们分析了什么问题，为什么重要
    \item \textbf{数据说明}（200字）：数据来源（仿真平台）、数据量、字段含义
    \item \textbf{分析方法}（200字）：用了 Hive SQL + Spark + Python，做了哪些分析
    \item \textbf{分析结果}（核心，配图表）：每张图配 2-3 句解读
    \item \textbf{结论建议}（200字）：根据分析给出业务建议
\end{enumerate}
\end{knowledgebox}

\begin{tipbox}[常见扣分点]
\begin{itemize}
    \item 只有数字没有图表
    \item 图表没有标题和坐标轴标签
    \item 只展示数据不做解读
    \item 没有结论建议
    \item 报告排版混乱
\end{itemize}
\end{tipbox}

\section{附录：SQL 速查表}
%========================================

{\footnotesize
\begin{longtable}{lll}
\toprule
\textbf{功能} & \textbf{SQL} & \textbf{说明} \\
\midrule
查询全部 & SELECT * FROM t & 查看所有数据 \\
条件查询 & SELECT * FROM t WHERE x=1 & 过滤行 \\
分组计数 & SELECT x, COUNT(*) FROM t GROUP BY x & 按x分组计数 \\
排序 & SELECT * FROM t ORDER BY x DESC & 降序排列 \\
限制条数 & SELECT * FROM t LIMIT 10 & 只取前10条 \\
去重 & SELECT DISTINCT x FROM t & 去除重复值 \\
聚合函数 & COUNT/SUM/AVG/MAX/MIN & 统计函数 \\
条件聚合 & SUM(CASE WHEN x THEN 1 ELSE 0 END) & 条件计数 \\
连接 & SELECT * FROM t1 JOIN t2 ON t1.id=t2.id & 多表连接 \\
子查询 & SELECT * FROM (SELECT ...) & 嵌套查询 \\
\bottomrule
\end{longtable}
}

\end{document}
